% ============================================================================
% 第八节：统一框架与开放问题
% ============================================================================

\section{统一框架与开放问题}\label{sec:unified}

基于前述各节，我们提出计数几何的统一递归框架，
并讨论若干开放问题与猜想。

\subsection{统一框架的表述}

\begin{conjecture}[递归统一猜想]\label{conj:unified}
\textbf{所有"自然"的计数不变量——
包括GW、DT、PT、Hurwitz数、Weil--Petersson体积、
矩阵模型自由能等——
都是某个谱曲线上CEO拓扑递归的输出。}
\end{conjecture}

这一猜想的精确化需要明确"自然"的含义。
基于现有证据，我们提出以下判据：

\begin{definition}[递归不变量的判据]\label{crit:recursive}
计数不变量$\{N_{g,n}\}$是"递归的"，若满足：
\begin{enumerate}[label=(\roman*)]
\item \textbf{切割公理}：$N_{g,n}$满足将$(g,n)$曲面切割为低复杂度曲面的递归；
\item \textbf{对称性}：$N_{g,n}$关于标记点对称；
\item \textbf{生成函数结构}：$F_{g,n} = \sum N_{g,n} \prod t_i$满足可积系统；
\item \textbf{模形式性}：在某些情形下，$N_{g,n}$与模形式相关。
\end{enumerate}
\end{definition}

\subsection{统一框架的技术基础}

统一框架的技术基础是以下三个等价性：

\begin{theorem}[三重等价性]\label{thm:triple}
以下三类对象在适当意义下等价：
\begin{enumerate}[label=(\arabic*)]
\item \textbf{CEO拓扑递归}：谱曲线$\mathcal{S}$上的多微分$\omega_{g,n}$；
\item \textbf{Givental量子化}：半单Frobenius流形上的$R$-矩阵作用；
\item \textbf{可积系统$\tau$-函数}：KP/Toda族中的特定$\tau$-函数。
\end{enumerate}
\end{theorem}

这一等价性由Dubrovin--Givental、Eynard--Orantin、Mulase--Zhang等人建立，
是计数几何统一性的技术核心。

\subsection{开放问题}

\begin{problem}[高亏格镜像对称的完全证明]\label{prob:mirror}
虽然亏格0镜像定理已完全证明，
高亏格情形（BCOV递归的严格化）仍是开放问题。
特别是：如何将BCOV方程的"边界条件"（holomorphic limit）
与CEO递归的初始数据精确对应？
\end{problem}

\begin{problem}[正特征计数几何]\label{prob:poschar}
GW不变量在正特征域上面临困难（形变论失效）。
DT/PT不变量在正特征下是否仍满足GW/DT/PT对应？
这一问题与SYLVA框架的"代数基础"相关。
\end{problem}

\begin{problem}[非半单Frobenius流形]\label{prob:nonsimplex}
Givental形式化要求量子上同调半单。
对非半单情形（如Calabi--Yau流形），
如何推广$R$-矩阵与CEO递归？
Teleman对非半单情形的工作提供了部分答案，但完整理论仍待建立。
\end{problem}

\begin{problem}[高维计数几何]\label{prob:highdim}
现有理论主要针对三维Calabi--Yau。
四维及更高维的计数几何（如Donaldson--Thomas理论的高维推广）
是否仍满足递归结构？
这与物理中的M-理论、F-理论相关。
\end{problem}

\begin{problem}[递归与物理对偶]\label{prob:physics}
CEO递归在物理上对应什么？
已有证据表明它与topological string的B-模型、
matrix model的大$N$极限相关，
但完整的物理诠释仍不清楚。
特别是：递归核$K_\alpha$的物理意义是什么？
\end{problem}

\subsection{可检验的猜想}

\begin{conjecture}[全亏格DT/PT对应的递归化]\label{conj:dtprec}
对任意光滑射影三维簇$X$，
存在谱曲线$\mathcal{S}_X$（由$X$的Gromov--Witten理论构造），
使得DT和PT不变量都是$\mathcal{S}_X$上CEO递归的输出。
\end{conjecture}

这一猜想在$X$为toric Calabi--Yau时已被验证（通过topological vertex），
对一般$X$仍是开放问题。

\begin{conjecture}[Catalan--Hurwitz--GW统一]\label{conj:CHG}
存在"普适谱曲线"$\mathcal{S}_{\mathrm{univ}}$，
使得Catalan数、Hurwitz数、$\CP^1$的GW不变量
都是$\mathcal{S}_{\mathrm{univ}}$上CEO递归在不同极限下的输出。
\end{conjecture}

这一猜想若成立，将揭示组合学（Catalan）、
复几何（Hurwitz）、代数几何（GW）的深层统一。

\subsection{与SYLVA框架的联系}

在SYLVA统一框架的语境下，
计数几何的递归统一性具有特殊意义：

\begin{remark}[SYLVA视角]\label{rmk:sylva}
SYLVA框架的"层级涌现"原理预测：
\textbf{不同层级的数学结构共享相同的生成机制}。
计数几何中GW/DT/PT/Hurwitz/Weil--Petersson的递归统一性
正是这一原理在数学中的体现——
它们都是"几何计数"这一深层结构在不同模空间上的投影。

CEO递归的核$K_\alpha$可视为这一深层结构的"生成元"，
而各类不变量是其"表示"。
这一视角为SYLVA框架的数学基础提供了严格支撑。
\end{remark}

\begin{remark}[可证伪性]\label{rmk:falsifiable}
统一框架的预测是可证伪的：
若发现某个"自然"计数不变量\textbf{不}满足CEO递归，
则框架需要修正。
目前所有已研究的自然计数不变量都满足递归，
但这一"经验规律"的严格化仍是挑战。
\end{remark}
